\import{order/hartogs.tex}

\section{Cardinal representatives}

\begin{definition}\label{cardinal}
    $\kappa$ is a cardinal iff $\kappa$ is an initial ordinal.
\end{definition}

\begin{proposition}\label{cardinal_is_ordinal}
    Let $\kappa$ be a cardinal.
    Then $\kappa$ is an ordinal.
\end{proposition}
\begin{proof}
    $\kappa$ is an initial ordinal by \obv{cardinal}.
    Follows by \obv{initial_ordinal}.
\end{proof}

For a set $X$, it is enough to search below its Hartogs ordinal for an
ordinal equinumerous with $X$.  This is an inhabited set-sized family, not a
selection from the proper class of ordinals.

\begin{definition}\label{cardinal_candidates}
    $\cardinalCandidates{X}
        =\{\alpha\in\hartogs{X}\mid
            \text{$\alpha$ is an ordinal and $\alpha\approx X$}\}$.
\end{definition}

\begin{definition}\label{card}
    $\card{X}=\inters{\cardinalCandidates{X}}$.
\end{definition}

\begin{proposition}\label{cardinal_candidates_inhabited}
    $\cardinalCandidates{X}$ is inhabited.
\end{proposition}
\begin{proof}
    Take $R$ such that $R$ is a strict well order on $X$
        by \obv{well_ordering_theorem}.
    Let $\alpha=\orderType{X}{R}$.
    $\alpha$ is an ordinal by \obv{order_type_is_ordinal}.
    $\alpha\in\hartogs{X}$
        by \auto{order_type_belongs_to_hartogs,subseteq_refl}.
    $\orderIsomorphism{\orderCollapse{X}{R}}
        {X}{R}{\alpha}{\memrel{\alpha}}$
        by \auto{order_collapse_is_isomorphism}.
    Hence $\orderCollapse{X}{R}\in\bijections{X}{\alpha}$
        by \obv{order_isomorphism}.
    Hence $X\approx\alpha$ by \obv{equinum}.
    Hence $\alpha\approx X$ by \obv{equinum_sym}.
    Thus $\alpha\in\cardinalCandidates{X}$
        by \obv{cardinal_candidates}.
\end{proof}

\begin{proposition}\label{cardinal_candidates_are_ordinals}
    Every element of $\cardinalCandidates{X}$ is an ordinal.
\end{proposition}
\begin{proof}
    Follows by \obv{cardinal_candidates}.
\end{proof}

\begin{proposition}\label{card_belongs_to_candidates}
    $\card{X}\in\cardinalCandidates{X}$.
\end{proposition}
\begin{proof}
    Follows by
        \auto{card,inters_of_ordinals_elem,cardinal_candidates_inhabited,cardinal_candidates_are_ordinals}.
\end{proof}

\begin{proposition}\label{card_is_ordinal}
    $\card{X}$ is an ordinal.
\end{proposition}
\begin{proof}
    Follows by \auto{card_belongs_to_candidates,cardinal_candidates}.
\end{proof}

\begin{proposition}\label{card_equinumerous}
    $\card{X}\approx X$.
\end{proposition}
\begin{proof}
    Follows by \auto{card_belongs_to_candidates,cardinal_candidates}.
\end{proof}

\begin{corollary}\label{equinumerous_card}
    $X\approx\card{X}$.
\end{corollary}
\begin{proof}
    Follows by \obv{card_equinumerous,equinum_sym}.
\end{proof}

\begin{theorem}\label{card_is_initial_ordinal}
    $\card{X}$ is an initial ordinal.
\end{theorem}
\begin{proof}
    $\card{X}$ is an ordinal by \obv{card_is_ordinal}.
    Show that for every $\beta\in\card{X}$ we have
        $\lnot(\beta\approx\card{X})$.
    \begin{subproof}
        Fix $\beta\in\card{X}$.
        $\beta$ is an ordinal by \obv{elem_of_ordinal_is_ordinal}.
        Suppose $\beta\approx\card{X}$.
        Hence $\beta\approx X$
            by \auto{card_equinumerous,equinum_tran}.
        $\card{X}\in\hartogs{X}$
            by \auto{card_belongs_to_candidates,cardinal_candidates}.
        $\hartogs{X}$ is an ordinal by \obv{hartogs_is_ordinal}.
        Hence $\beta\in\hartogs{X}$ by \obv{ordinal_transitivity}.
        Hence $\beta\in\cardinalCandidates{X}$
            by \obv{cardinal_candidates}.
        $\card{X}\subseteq\beta$
            by
            \auto{card,inters_subseteq_elem,cardinal_candidates_inhabited}.
        Hence $\beta\in\beta$ by \obv{subseteq}.
        Contradiction by \obv{in_irrefl}.
    \end{subproof}
    Follows by \obv{initial_ordinal}.
\end{proof}

\begin{theorem}\label{card_is_cardinal}
    $\card{X}$ is a cardinal.
\end{theorem}
\begin{proof}
    Follows by \obv{cardinal,card_is_initial_ordinal}.
\end{proof}

\begin{theorem}\label{card_least_equinumerous_ordinal}
    Let $\alpha$ be an ordinal.
    Suppose $\alpha\approx X$.
    Then $\card{X}\subseteq\alpha$.
\end{theorem}
\begin{proof}
    $\alpha\approx\card{X}$
        by \auto{equinumerous_card,equinum_tran}.
    Follows by
        \obv{card_is_initial_ordinal,initial_ordinal_le_equinumerous_ordinal}.
\end{proof}

\begin{theorem}\label{card_eq_iff_equinumerous}
    $\card{X}=\card{Y}$ iff $X\approx Y$.
\end{theorem}
\begin{proof}
    Show that if $\card{X}=\card{Y}$, then $X\approx Y$.
    \begin{subproof}
        Assume $\card{X}=\card{Y}$.
        $X\approx\card{X}$ by \obv{equinumerous_card}.
        $\card{Y}\approx Y$ by \obv{card_equinumerous}.
        Follows by \auto{equinum_tran}.
    \end{subproof}
    Show that if $X\approx Y$, then $\card{X}=\card{Y}$.
    \begin{subproof}
        Assume $X\approx Y$.
        $\card{X}\approx X$ by \obv{card_equinumerous}.
        $Y\approx\card{Y}$ by \obv{equinumerous_card}.
        Hence $\card{X}\approx\card{Y}$ by \obv{equinum_tran}.
        $\card{X}$ and $\card{Y}$ are initial ordinals
            by \obv{card_is_initial_ordinal}.
        Follows by \obv{equinumerous_initial_ordinals_are_equal}.
    \end{subproof}
    Follows by \auto.
\end{proof}

\begin{proposition}\label{card_of_cardinal}
    Let $\kappa$ be a cardinal.
    Then $\card{\kappa}=\kappa$.
\end{proposition}
\begin{proof}
    $\card{\kappa}$ is an initial ordinal
        by \obv{card_is_initial_ordinal}.
    $\card{\kappa}\approx\kappa$ by \obv{card_equinumerous}.
    Follows by \obv{equinumerous_initial_ordinals_are_equal}.
\end{proof}

\begin{proposition}\label{card_fixed_iff_cardinal}
    $\card{X}=X$ iff $X$ is a cardinal.
\end{proposition}
\begin{proof}
    Follows by \auto{card_is_cardinal,card_of_cardinal}.
\end{proof}
